\chapter{테스트 관리}

\begin{summarybox}{한눈에}
본 기관이 만든 \textbf{기타 테스트}만 보입니다. 본사가 만든 진단 / 챕터 시험은 사이드바에서 자동 필터되어 보이지 않아요. 따라서 종류 탭(전체/진단/챕터/기타)도 숨겨지고, 화면이 더 단순합니다.
\end{summarybox}

\kfShot{09-tests}{테스트 관리 — 본 기관이 만든 시험만}

\section*{화면 구조}

\begin{screenbox}{시험 목록 + 신규 생성}
\noindent
\begin{tabular}{@{}p{3cm}p{2cm}p{2cm}p{2cm}p{2cm}p{2cm}@{}}
\toprule
\textbf{시험명} & \textbf{레벨} & \textbf{문항} & \textbf{제출 수} & \textbf{시험일} & \textbf{액션} \\
\midrule
초3 1학기 중간 & 소쉬르 3 & 20 & 8 & 2026-04-15 & \inacttab{편집} \inacttab{통계} \\
중1 단원 평가 1 & 러셀 1 & 30 & 12 & 2026-04-22 & \inacttab{편집} \inacttab{통계} \\
\bottomrule
\end{tabular}

\medskip
\activetab{+ 시험 추가}
\end{screenbox}

\section*{여기서 할 수 있는 일}
\begin{itemize}[leftmargin=1.2em]
  \item \textbf{시험 추가} — 비주얼 에디터로 이동
  \item \textbf{편집} — 기존 시험 문항·정답·해설 수정
  \item \textbf{통계} — 학생별 응시 결과·정답률·평균 점수
  \item 시험 PDF 미리보기 + 정답표 별도 출력
  \item 검색·필터 (제목·레벨·시리즈)
\end{itemize}

\section*{자주 쓰는 흐름}
\begin{enumerate}
  \item ``+ 시험 추가'' → 제목·레벨 입력 → 비주얼 에디터 진입
  \item 문항 작성 또는 PDF 업로드 후 OCR → 자동 변환
  \item 학습 계획표 ``테스트'' 열에 배정
  \item 학생 응시 후 ``통계'' 에서 결과 확인
\end{enumerate}

\begin{tipbox}{알아두면 좋은 점}
\begin{itemize}[leftmargin=1.2em]
  \item 본 기관이 만든 시험만 보입니다 — 본사 PUBLIC 시험·진단·챕터 모두 자동 필터
  \item 종류 탭(전체/진단/챕터/기타) 자체가 \textbf{숨김} — ORG\_ADMIN 이 ``기타 테스트'' 만 다룬다는 정책
  \item 시험은 본 기관 학생에게만 배정 가능
\end{itemize}
\end{tipbox}
